\section{TxAuthEnvelope: a 17-field signed transaction}
\label{sec:tx-auth-envelope}

\subsection{What it replaces}

Every classical Ethereum-family transaction lineage --- legacy raw RLP,
EIP-2930 access list, EIP-1559 dynamic fee, EIP-4844 blob --- shares the same
authorisation primitive: ECDSA over keccak256 of an RLP-encoded payload. The
payload format changes across tx types but the authorisation primitive is
fixed. Both are classical: ECDSA falls to Shor, keccak is fine but the
domain-separation gaps in RLP encoding admit cross-version replays
\cite{eip2718,eip7932}.

A strict-PQ chain needs a transaction-authorisation primitive that (1) signs
under ML-DSA-65 or another FIPS-204 scheme, (2) binds every security-relevant
field via TupleHash256 so a single byte flip on any field breaks the
signature, (3) carries the profile identity on the wire so cross-profile
replays are impossible, and (4) survives lazy-deserialiser bugs by refusing to
admit a transaction whose decoded fields disagree with its decoded digest.

\texttt{TxAuthEnvelope} (\texttt{luxfi/consensus@v1.23.6},
\texttt{protocol/auth/tx\_envelope.go:108}) is that primitive.

\subsection{The 17 fields}

\begin{table}[H]
\centering
\caption{TxAuthEnvelope field list. Source:
\texttt{protocol/auth/tx\_envelope.go:108}.}
\label{tab:tx-auth-fields}
\begin{tabular}{llp{7cm}}
\toprule
\# & Field & Role \\
\midrule
1  & \texttt{Version}          & envelope version (uint16). $v=1$ on strict-PQ. \\
2  & \texttt{ProfileID}        & profile under which signed. \\
3  & \texttt{ChainID}          & destination chain id. \\
4  & \texttt{NetworkID}        & destination network id (mainnet / testnet). \\
5  & \texttt{AccountID}        & 48-byte signer account identifier. \\
6  & \texttt{Nonce}            & uint64 monotonic per-account nonce. \\
7  & \texttt{ExpiryHeight}     & uint64 chain height after which envelope is invalid. \\
8  & \texttt{WalletSchemeID}   & signature scheme (matches \texttt{profile.WalletSchemeID}). \\
9  & \texttt{HashSuiteID}      & hash suite (matches \texttt{profile.HashSuiteID}). \\
10 & \texttt{FeePayer}         & 48-byte account paying gas. \\
11 & \texttt{GasLimit}         & uint64 gas allowance. \\
12 & \texttt{MaxFee}           & 32-byte EIP-1559-style fee cap. \\
13 & \texttt{CallRoot}         & 32-byte hash of the call payload. \\
14 & \texttt{AccessListRoot}   & 32-byte hash of the access list. \\
15 & \texttt{ZIdentityRoot}    & 32-byte Z-Chain identity root binding. \\
16 & \texttt{AccountStateRoot} & 32-byte expected account state root at submit. \\
17 & \texttt{PublicKeyRef}     & 32-byte SHA3-256 of the signer's full ML-DSA pubkey. \\
\bottomrule
\end{tabular}
\end{table}

The 17 fields are the security-relevant inputs to the signing transcript. The
sole non-transcript field is \texttt{Signature}, which carries the
3309-byte ML-DSA-65 signature itself (variable-length under other schemes).

The choice of 17 fields is the smallest set that closes \emph{every} replay
class the protocol team enumerated in HIP-0085 §3:

\begin{itemize}[leftmargin=2em,nosep]
  \item \textbf{Cross-chain replay} is closed by \texttt{ChainID} and
    \texttt{NetworkID}.
  \item \textbf{Cross-profile replay} is closed by \texttt{ProfileID}.
  \item \textbf{Cross-version replay} is closed by \texttt{Version}.
  \item \textbf{Cross-account replay} is closed by \texttt{AccountID} and
    \texttt{PublicKeyRef}.
  \item \textbf{Cross-fee-payer replay} is closed by \texttt{FeePayer},
    \texttt{GasLimit}, \texttt{MaxFee}.
  \item \textbf{Stale state replay} is closed by \texttt{ExpiryHeight},
    \texttt{AccountStateRoot}, and \texttt{ZIdentityRoot}.
  \item \textbf{Cross-permission-set replay} is closed by
    \texttt{AccessListRoot}.
  \item \textbf{Same-account same-state different-call replay} is closed by
    \texttt{CallRoot}.
  \item \textbf{Cross-scheme replay} is closed by \texttt{WalletSchemeID}
    and \texttt{HashSuiteID}.
\end{itemize}

Dropping any field would re-open a replay class. The 17-field minimum is
audit-driven.

\subsection{SigningDigest}

\texttt{SigningDigest} (\texttt{tx\_envelope.go:148}) is what the wallet signs:

\begin{lstlisting}[language=Go,caption={\texttt{SigningDigest} construction.
Source: \texttt{protocol/auth/tx\_envelope.go:148}. Hash family is cSHAKE256
under \texttt{LUX\_TX\_AUTH\_V1}, independent of \texttt{HashSuiteID} which
is bound as data into the transcript.}]
func (e *TxAuthEnvelope) SigningDigest() [48]byte {
    parts := [][]byte{
        []byte(txAuthProtocolTag),       // "Lux/TxAuthEnvelope/v1"
        u16BE(e.Version),
        {byte(e.ProfileID)},
        u32BE(e.ChainID),
        u32BE(e.NetworkID),
        e.AccountID[:],
        u64BE(e.Nonce),
        u64BE(e.ExpiryHeight),
        {byte(e.WalletSchemeID)},
        {byte(e.HashSuiteID)},
        e.FeePayer[:],
        u64BE(e.GasLimit),
        e.MaxFee[:],
        e.CallRoot[:],
        e.AccessListRoot[:],
        e.ZIdentityRoot[:],
        e.AccountStateRoot[:],
        e.PublicKeyRef[:],
    }
    return tupleHash48(parts, txAuthSigningCustomization)
}
\end{lstlisting}

The hash family is cSHAKE256, output truncated to 48 bytes; the customization
string is \texttt{LUX\_TX\_AUTH\_V1}. \texttt{HashSuiteID} is bound \emph{as
data} into the transcript --- it is not the hash family of the transcript
itself. This means a profile change to \texttt{HashSuiteID} bumps the
customization (\texttt{LUX\_TX\_AUTH\_V1} $\to$ \texttt{V2}) at the same time
it changes the data byte. Both changes are observable in the digest.

\subsection{Verification gate}

\texttt{VerifyTxAuthEnvelope} (\texttt{protocol/auth/tx\_envelope.go}) is the
profile-gated verifier. The order matters; we list it explicitly because
the order is the security-relevant claim:

\begin{enumerate}[leftmargin=2em,nosep]
  \item \textbf{Structural.} Envelope non-nil, profile non-nil,
    \texttt{Version} supported, every field bytes-of-the-right-width.
  \item \textbf{Profile/id match.} \texttt{envelope.ProfileID} equals
    \texttt{profile.ProfileID}.
  \item \textbf{Profile/scheme match.} \texttt{envelope.WalletSchemeID}
    equals \texttt{profile.WalletSchemeID}.
  \item \textbf{Profile/suite match.} \texttt{envelope.HashSuiteID} equals
    \texttt{profile.HashSuiteID}.
  \item \textbf{Profile typed-auth gate.}
    \texttt{profile.RequireTypedTxAuth} is \texttt{true}; if false this
    code path is not reached at all.
  \item \textbf{Z-Chain binding.} \texttt{ZIdentityRoot} matches the
    current Z-Chain identity root for this epoch (look-up against
    \texttt{ZRegistryRoots.IdentityRoot}, Section~\ref{sec:zchain-rollup}).
  \item \textbf{Account state binding.} \texttt{AccountStateRoot} matches
    the current account state root.
  \item \textbf{Account / pubkey binding.} The resolved pubkey hashes to
    \texttt{PublicKeyRef}, and the pubkey is on the active key list for
    \texttt{AccountID} as of \texttt{ZIdentityRoot}.
  \item \textbf{Expiry.} Current chain height $\le$ \texttt{ExpiryHeight}.
  \item \textbf{Nonce.} \texttt{envelope.Nonce} equals
    \texttt{account.Nonce} (no fast-forward, no replay).
  \item \textbf{Signature.} ML-DSA-65 verify of \texttt{Signature} over
    \texttt{SigningDigest()} under the resolved pubkey, FIPS 204 §5.2 ctx
    set to a per-profile constant.
\end{enumerate}

Steps 2--5 reject before the expensive signature verification runs. Step 6--8
reject before a stale-state replay can land. Step 11 is the only step that
runs the lattice signature primitive, which is the only PQ-expensive step.

\subsection{Comparison to EIP-7932}

EIP-7932 \cite{eip7932} proposes a signature registry for EVM transactions
that lets new schemes plug in at the transaction-type byte. The EIP is
structurally compatible with this paper's goals --- it would allow a chain to
admit ML-DSA-65 transactions alongside ECDSA transactions --- but it falls
short in two ways:

\begin{itemize}[leftmargin=2em,nosep]
  \item EIP-7932 does not standardise the \emph{transcript}. Each scheme
    declares its own signing digest. Two schemes sharing inputs can produce
    different digests because the scheme-byte appears in the wrapper but
    not in the inner transcript. TupleHash256 with a single customization
    closes that gap.
  \item EIP-7932 does not standardise the profile binding. A transaction
    signed under scheme A on chain X can be replayed as a scheme-A
    transaction on chain Y if both chains admit scheme A in their
    registry. \texttt{ProfileID} closes that gap.
\end{itemize}

Bullets aside, the two designs converge: a strict-PQ chain that wants to
expose a 7932-style registry to other chains can declare exactly one entry
(\texttt{TxAuthEnvelope}) and exactly one scheme (ML-DSA-65), and the
externally-observable behaviour is equivalent. The Lux production canonical
takes the simpler path: no in-chain registry, a single mandatory envelope.

\subsection{Wire codec}

\texttt{TxAuthEnvelope} serialises as a deterministic big-endian layout
(\texttt{tx\_envelope.go:178}):

\begin{table}[H]
\centering
\caption{TxAuthEnvelope wire layout, fixed-width except \texttt{signature}.}
\begin{tabular}{lll}
\toprule
Offset & Size & Field \\
\midrule
$0$    & $2$   & \texttt{version}            (uint16 BE) \\
$2$    & $1$   & \texttt{profile\_id}        (uint8) \\
$3$    & $4$   & \texttt{chain\_id}          (uint32 BE) \\
$7$    & $4$   & \texttt{network\_id}        (uint32 BE) \\
$11$   & $48$  & \texttt{account\_id} \\
$59$   & $8$   & \texttt{nonce}              (uint64 BE) \\
$67$   & $8$   & \texttt{expiry\_height}     (uint64 BE) \\
$75$   & $1$   & \texttt{wallet\_scheme\_id} (uint8) \\
$76$   & $1$   & \texttt{hash\_suite\_id}    (uint8) \\
$77$   & $48$  & \texttt{fee\_payer} \\
$125$  & $8$   & \texttt{gas\_limit}         (uint64 BE) \\
$133$  & $32$  & \texttt{max\_fee} \\
$165$  & $32$  & \texttt{call\_root} \\
$197$  & $32$  & \texttt{access\_list\_root} \\
$229$  & $32$  & \texttt{z\_identity\_root} \\
$261$  & $32$  & \texttt{account\_state\_root} \\
$293$  & $32$  & \texttt{public\_key\_ref} \\
$325$  & $4$   & \texttt{sig\_length}        (uint32 BE) \\
$329$  & $\le$ 3309 & \texttt{signature}     (ML-DSA-65) \\
\bottomrule
\end{tabular}
\end{table}

The 329-byte fixed header followed by a length-prefixed signature is the
entire wire format. There is no nested encoding, no varints, no RLP
ambiguity, no field-ordering choice. A naive decoder that consumes the bytes
in order produces the canonical envelope; \emph{any} ambiguity in decoding is
a fatal decoder bug, not a protocol question.
